% GENERATED FILE. Edit canonical lessons, then run npm run generate:books.
% canonical-source-hash: fnv1a64:ed8cdef851eae306
% canonical-lessons: RU-C533-shurup, RU-C533-ploskogubtsy, RU-C533-otkruchivat, RU-C533-peregorat, RU-C533-izolenta

\chapter{Small Repairs}
\label{ch:ru-small-repairs}
\hlchaptermodality{533}

\begin{chapteropening}
\textbf{By the end of this chapter:} \emph{I can talk about small repairs with screws, pliers and tape.}

Complete the last lesson of chapter 533: I can talk about small repairs with screws, pliers and tape.
\end{chapteropening}

\section[\texorpdfstring{\ru{шуруп} (shurúp)}{шуруп (shurúp)}]{\ru{шуруп} (shurúp) — a screw}

\label{lesson:RU-C533-shurup}



\begin{quote}
Before the new one: say the Russian for a power drill, then the Russian for a screwdriver.
\end{quote}

\subsection*{You'll want to know: \ru{шуруп}}
\textbf{\ru{шуруп}} (\emph{shurúp}) — \textquotedblleft{}a screw\textquotedblright{}.

\begin{quote}
\ru{Для} \ru{полки} \ru{нужны} \ru{четыре} \ru{шурупа}.
\end{quote}

\begin{quote}
\emph{Dlya pólki nuzhný chetýre shurúpa.}
\end{quote}

\begin{quote}
The shelf needs four screws.
\end{quote}

Masculine.

\begin{cousinweb}
\emph{\ru{шуруп}} was borrowed from German \emph{Schraube} 'a screw' (EtymDB).
\end{cousinweb}

\begin{grammarlens}[title={What you've built}]
One more word for telling, explaining and describing at B1.
\end{grammarlens}

\subsection*{Your turn}
\begin{itemize}
\raggedright
  \item \emph{Say it:} \emph{shurúp}
  \item \emph{Say it:} \emph{shurúp}, once more
  \item \emph{Say it:} a sentence of your own with \emph{shurúp}
  \item \emph{Recall:} say the Russian for a power drill, then the Russian for a screwdriver, then say \emph{shurúp} again
\end{itemize}

\begin{tcolorbox}[breakable,colback=teal!4,colframe=teal!35!black,title={Before you move on}]
What does \ru{шуруп} mean? (\textquotedblleft{}A screw\textquotedblright{}.) Say it once more.
\end{tcolorbox}
\section[\texorpdfstring{\ru{плоскогубцы} (ploskogúbtsy)}{плоскогубцы (ploskogúbtsy)}]{\ru{плоскогубцы} (ploskogúbtsy) — pliers}

\label{lesson:RU-C533-ploskogubtsy}



\begin{quote}
Before the new one: say the Russian for a screwdriver, then the Russian for a screw.
\end{quote}

\subsection*{You'll want to know: \ru{плоскогубцы}}
\textbf{\ru{плоскогубцы}} (\emph{ploskogúbtsy}) — \textquotedblleft{}pliers\textquotedblright{}.

\begin{quote}
\ru{Плоскогубцы} \ru{лежат} \ru{в} \ru{ящике} \ru{с} \ru{инструментами}.
\end{quote}

\begin{quote}
\emph{Ploskogúbtsy lezhát v yáshchike s instruméntami.}
\end{quote}

\begin{quote}
The pliers are in the toolbox.
\end{quote}

Plural only.

\begin{cousinweb}
\emph{\ru{плоскогубцы}} is a compound of \emph{\ru{плоский}} 'flat' and \emph{\ru{губа}} 'lip': pliers with flat jaws; EtymDB has no entry.
\end{cousinweb}

\begin{grammarlens}[title={What you've built}]
One more word for telling, explaining and describing at B1.
\end{grammarlens}

\subsection*{Your turn}
\begin{itemize}
\raggedright
  \item \emph{Say it:} \emph{ploskogúbtsy}
  \item \emph{Say it:} \emph{ploskogúbtsy}, once more
  \item \emph{Say it:} a sentence of your own with \emph{ploskogúbtsy}
  \item \emph{Recall:} say the Russian for a screwdriver, then the Russian for a screw, then say \emph{ploskogúbtsy} again
\end{itemize}

\begin{tcolorbox}[breakable,colback=teal!4,colframe=teal!35!black,title={Before you move on}]
What does \ru{плоскогубцы} mean? (\textquotedblleft{}Pliers\textquotedblright{}.) Say it once more.
\end{tcolorbox}
\section[\texorpdfstring{\ru{откручивать} (otkrúchivat)}{откручивать (otkrúchivat)}]{\ru{откручивать} (otkrúchivat) — to unscrew}

\label{lesson:RU-C533-otkruchivat}



\begin{quote}
Before the new one: say the Russian for a screw, then the Russian for pliers.
\end{quote}

\subsection*{You'll want to know: \ru{откручивать}}
\textbf{\ru{откручивать}} (\emph{otkrúchivat}) — \textquotedblleft{}to unscrew\textquotedblright{}.

\begin{quote}
\ru{Папа} \ru{откручивает} \ru{старую} \ru{лампочку}.
\end{quote}

\begin{quote}
\emph{Pápa otkrúchivayet stáruyu lámpochku.}
\end{quote}

\begin{quote}
Dad is unscrewing the old light bulb.
\end{quote}

Imperfective; perfective \ru{открутить} (\ru{откручу}).

\begin{cousinweb}
\emph{\ru{откручивать}} is built on \emph{\ru{от}-} 'off' and \emph{\ru{крутить}} 'to twist, to turn'; EtymDB has no entry.
\end{cousinweb}

\begin{grammarlens}[title={What you've built}]
One more word for telling, explaining and describing at B1.
\end{grammarlens}

\subsection*{Your turn}
\begin{itemize}
\raggedright
  \item \emph{Say it:} \emph{otkrúchivat}
  \item \emph{Say it:} \emph{otkrúchivat}, once more
  \item \emph{Say it:} a sentence of your own with \emph{otkrúchivat}
  \item \emph{Recall:} say the Russian for a screw, then the Russian for pliers, then say \emph{otkrúchivat} again
\end{itemize}

\begin{tcolorbox}[breakable,colback=teal!4,colframe=teal!35!black,title={Before you move on}]
What does \ru{откручивать} mean? (\textquotedblleft{}To unscrew\textquotedblright{}.) Say it once more.
\end{tcolorbox}
\section[\texorpdfstring{\ru{перегорать} (peregorát)}{перегорать (peregorát)}]{\ru{перегорать} (peregorát) — to burn out (of a bulb)}

\label{lesson:RU-C533-peregorat}



\begin{quote}
Before the new one: say the Russian for pliers, then the Russian for to unscrew.
\end{quote}

\subsection*{You'll want to know: \ru{перегорать}}
\textbf{\ru{перегорать}} (\emph{peregorát}) — \textquotedblleft{}to burn out (of a bulb)\textquotedblright{}.

\begin{quote}
\ru{На} \ru{кухне} \ru{часто} \ru{перегорают} \ru{лампочки}.
\end{quote}

\begin{quote}
\emph{Na kúkhne chásto peregoráyut lámpochki.}
\end{quote}

\begin{quote}
Light bulbs often blow in the kitchen.
\end{quote}

Imperfective; perfective \ru{перегореть}: \emph{\ru{лампочка} \ru{перегорела}}.

\begin{cousinweb}
\emph{\ru{перегорать}} is built on \emph{\ru{пере}-} and \emph{\ru{гореть}} 'to burn'; EtymDB has no entry.
\end{cousinweb}

\begin{grammarlens}[title={What you've built}]
One more word for telling, explaining and describing at B1.
\end{grammarlens}

\subsection*{Your turn}
\begin{itemize}
\raggedright
  \item \emph{Say it:} \emph{peregorát}
  \item \emph{Say it:} \emph{peregorát}, once more
  \item \emph{Say it:} a sentence of your own with \emph{peregorát}
  \item \emph{Recall:} say the Russian for pliers, then the Russian for to unscrew, then say \emph{peregorát} again
\end{itemize}

\begin{tcolorbox}[breakable,colback=teal!4,colframe=teal!35!black,title={Before you move on}]
What does \ru{перегорать} mean? (\textquotedblleft{}To burn out (of a bulb)\textquotedblright{}.) Say it once more.
\end{tcolorbox}
\section[\texorpdfstring{\ru{изолента} (izolénta)}{изолента (izolénta)}]{\ru{изолента} (izolénta) — insulating tape}

\label{lesson:RU-C533-izolenta}



\begin{quote}
Before the new one: say the Russian for to unscrew, then the Russian for to burn out.
\end{quote}

\subsection*{You'll want to know: \ru{изолента}}
\textbf{\ru{изолента}} (\emph{izolénta}) — \textquotedblleft{}insulating tape\textquotedblright{}.

\begin{quote}
\ru{Он} \ru{обмотал} \ru{провод} \ru{изолентой}.
\end{quote}

\begin{quote}
\emph{On obmotál próvod izoléntoy.}
\end{quote}

\begin{quote}
He wrapped the wire in insulating tape.
\end{quote}

Feminine. Short for \emph{\ru{изоляционная} \ru{лента}}.

\begin{cousinweb}
\emph{\ru{изолента}} is a clipped compound of \emph{\ru{изоляционная}} 'insulating' and \emph{\ru{лента}} 'a ribbon, a tape'; EtymDB has no entry.
\end{cousinweb}

\begin{grammarlens}[title={What you've built}]
One more word for telling, explaining and describing at B1.
\end{grammarlens}

\subsection*{Your turn}
\begin{itemize}
\raggedright
  \item \emph{Say it:} \emph{izolénta}
  \item \emph{Say it:} \emph{izolénta}, once more
  \item \emph{Say it:} a sentence of your own with \emph{izolénta}
  \item \emph{Recall:} say the Russian for to unscrew, then the Russian for to burn out, then say \emph{izolénta} again
\end{itemize}

\begin{tcolorbox}[breakable,colback=teal!4,colframe=teal!35!black,title={Before you move on}]
What does \ru{изолента} mean? (\textquotedblleft{}Insulating tape\textquotedblright{}.) Say it once more.
\end{tcolorbox}
